\documentclass[11pt,a4paper]{article}
\usepackage[utf8]{inputenc}
\usepackage[T1]{fontenc}
\usepackage[margin=1in]{geometry}
\usepackage{xcolor}
\usepackage{hyperref}
\usepackage{lmodern}
\definecolor{primarycolor}{RGB}{33,37,41}
\definecolor{secondarycolor}{RGB}{0,102,204}
\hypersetup{colorlinks=true,linkcolor=secondarycolor,urlcolor=secondarycolor,citecolor=secondarycolor}
\pagestyle{empty}
\setlength{\parindent}{0pt}
\setlength{\parskip}{0.8em}
\begin{document}
{\Large\bfseries\color{primarycolor} Marta Kowal-Brandt}\\
{\small \href{mailto:marta.kowalbrandt@example.com}{marta.kowalbrandt@example.com} | +48 555 303 303 | \href{https://gitlab.com/mkb-example}{gitlab.com/mkb-example}}

Corvale Mobility\\
\textbf{Application: Cloud Infrastructure / DevOps Contractor -- Fleet Connectivity}

Dear Corvale Mobility hiring team,

At Odrava Health Software, I own GitLab CI pipelines for three product teams supporting approximately 25 developers. My work has increasingly focused on deployment automation, ephemeral environments and build tooling, and I am applying for your Cloud Infrastructure / DevOps Contractor role as I move into DevOps / platform engineering full time. The part of your fleet-connectivity scope where I can contribute most directly is helping backend developers deploy services that are easier to monitor and run.

I reduced average pipeline time from 38 to 14 minutes through parallelisation, dependency caching and changes to end-to-end test execution, while monthly CI runner costs fell by approximately 30\%. I also built merge-request environments that deploy the application, PostgreSQL and seeded data through GitLab CI and Helm, then clean themselves up automatically. These environments let QA and product review features before merge. Alongside the implementation, I wrote the internal CI documentation used in onboarding, making the workflows easier for others to understand and use.

My infrastructure experience combines workplace test environments with personal projects. At work, I deploy within namespaces on a Kubernetes cluster operated by the platform team. In my homelab, I built a three-node k3s cluster with Argo CD, Prometheus, Grafana, Loki, Traefik and cert-manager for TLS, and provisioned a personal Hetzner VPS with Terraform. My Terraform experience is personal rather than production-based, and my AWS experience is limited to learning EC2 and S3 basics. This is a transition from test infrastructure, rather than a claim to have already led production cloud architecture.

Corvale's fleet-connectivity platform brings deployment, observability and reliability together from the start. I would bring measurable CI/CD improvements, hands-on environment automation and clear documentation to that work. I would welcome a discussion about how this experience could contribute to your team and where it fits the engagement's scope.

Kind regards,\\
Marta Kowal-Brandt
\end{document}